\documentclass[11pt]{article}

\usepackage[utf8]{inputenc}
\usepackage[T1]{fontenc}
\usepackage{lmodern}
\usepackage[margin=1in]{geometry}
\usepackage{amsmath,amssymb}
\usepackage{booktabs}
\usepackage{multirow}
\usepackage{graphicx}
\usepackage{microtype}
\usepackage{algorithm}
\usepackage{algpseudocode}
\usepackage[round]{natbib}
\usepackage[colorlinks=true,linkcolor=blue,citecolor=blue,urlcolor=blue]{hyperref}

\newcommand{\logit}{\operatorname{logit}}
\newcommand{\sigm}{\sigma}
\newcommand{\RR}{\mathbb{R}}
\newcommand{\EE}{\mathbb{E}}

\title{\textbf{Pattern-Expert Additive Models:\\ Interpretable Classifiers that Win on\\ Localized High-Order Interactions}}
\author{Rupesh Bharambe\\ \small Independent Researcher \\ \small \texttt{rupeshbharambe2004@gmail.com}}
\date{\today}

\begin{document}
\maketitle

\begin{abstract}
Generalized additive models with pairwise interactions (GA2M), such as the
Explainable Boosting Machine (EBM), are the accuracy leaders among intrinsically
interpretable models for tabular classification, but they cannot represent
interactions of order three or higher. Black-box gradient boosting captures such
interactions yet requires post-hoc, approximate explanations. We study the
\emph{Pattern-Expert Additive Model} (PEAM), which partitions features into
groups, fits one calibrated, arbitrarily flexible expert per group, and combines
the experts additively in log-odds space. Each prediction decomposes
\emph{exactly} into signed per-group reason codes, with no post-hoc
approximation; interactions of any order are captured \emph{within} a group,
while the model stays additive \emph{across} groups. We show --- honestly --- that
PEAM is not a general-purpose state of the art: on linear, pairwise, saturated,
or cross-group-interaction data it is matched or beaten by simpler models.
Instead it occupies a precise and previously uncharted point on the
accuracy--interpretability frontier. Through a controlled synthetic study that
sweeps interaction order, and experiments on four public datasets, we
characterize exactly when PEAM wins: when the predictive signal consists of
\emph{high-order interactions localized within a known feature grouping}. There,
PEAM is the best interpretable model and can beat strong black boxes. On the
MADELON benchmark PEAM attains ROC-AUC $0.90$ versus $0.80$ for EBM and $0.86$
for full-feature XGBoost; on HIGGS, whose high-order signal is cross-object
(delocalized), PEAM correctly fails --- a contrast that isolates locality, not
merely interaction order, as the operative condition. We also show that which
grouping method succeeds is structure-dependent, and that neither correlation
nor univariate mutual information is universally adequate.
\end{abstract}

\section{Introduction}

Tabular classifiers face a persistent tension between accuracy and
interpretability. At one end, generalized additive models
\citep{hastie1990gam} and their modern incarnations --- in particular GA2M and
the Explainable Boosting Machine (EBM) \citep{lou2013ga2m,caruana2015intelligible,nori2019interpretml}
--- are intrinsically interpretable: every prediction is a sum of univariate shape
functions and (for GA2M) pairwise interaction terms, each of which can be plotted
and audited. At the other end, gradient-boosted tree ensembles
\citep{chen2016xgboost,ke2017lightgbm} achieve higher accuracy by modelling
interactions of arbitrary order, but they are black boxes whose explanations must
be reconstructed post hoc with methods such as SHAP \citep{lundberg2017shap} or
LIME \citep{ribeiro2016lime}, which are approximate and can be unfaithful to the
underlying model.

A structural limitation sits at the heart of this trade-off: GA2M, by
construction, is additive in main effects and \emph{pairwise} interactions. It
\emph{cannot} represent a genuine third- or higher-order interaction. When the
label depends on such an interaction, GA2M is misspecified and a black box wins
--- but then interpretability is lost. This motivates a simple question:
\emph{is there an intrinsically interpretable model that captures high-order
interactions exactly where they occur, while producing faithful attributions?}

We study one answer, the \textbf{Pattern-Expert Additive Model (PEAM)}. PEAM
partitions the features into $K$ groups, fits a single calibrated but otherwise
\emph{arbitrarily flexible} expert on each group, and combines the experts in
log-odds (evidence) space:
\begin{equation}
\logit\, p(x) \;=\; b_0 \;+\; \sum_{k=1}^{K} \beta_k\,\big[\logit e_k(x_{G_k}) - \logit \pi\big],
\label{eq:peam-intro}
\end{equation}
where $e_k$ is the group-$k$ expert and $\pi$ the base rate. Because the
combination is additive, every prediction decomposes \emph{exactly} into signed
per-group reason codes $r_k = \beta_k[\logit e_k - \logit \pi]$ that sum to the
decision log-odds --- no post-hoc approximation. Because each expert is
unconstrained, interactions of \emph{any} order are captured \emph{within} a
group; because the experts are combined additively, the model assumes no
interaction \emph{across} groups.

Our central message is deliberately honest and double-edged. PEAM is
\textbf{not} a general-purpose state of the art. On data that is linear,
dominated by pairwise interactions, near-separable, or whose interactions span
groups, PEAM is matched or beaten by logistic regression, by EBM, or by a black
box. What PEAM offers is a \emph{new, precisely characterizable} operating
regime: when the predictive signal is composed of \emph{high-order interactions
that are localized within a known grouping}, PEAM's inductive bias is exactly
right, and it becomes the best interpretable model --- at times beating even
black boxes, because its structure additionally performs implicit structured
feature selection. We make this precise.

\paragraph{Contributions.}
\begin{itemize}
\item We formalize PEAM (Section~\ref{sec:method}): a group-additive classifier
with flexible, calibrated per-group experts and exact additive log-odds reason
codes, and we state the inductive bias it encodes.
\item We give a specification analysis (Section~\ref{sec:when}) identifying the
regime in which PEAM is well-specified and GA2M is not, and the regime
(cross-group interactions) in which PEAM is misspecified.
\item We verify the analysis with a controlled synthetic study
(Section~\ref{sec:synthetic}) that sweeps interaction order and sample size: at
order $\geq 3$ PEAM is the single best model, beating EBM and strong, tuned
black boxes from $n{=}2$k to $n{=}80$k.
\item We map where the regime occurs in public data
(Section~\ref{sec:real}). On MADELON \citep{guyon2003madelon} PEAM is the best
interpretable model and beats full-feature black boxes; on HIGGS
\citep{baldi2014higgs} its high-order-but-delocalized signal makes PEAM fail, as
predicted; on saturated / low-order data (satimage, gas-drift) standard methods
win.
\item We show that the adequacy of a grouping is structure-dependent
(Section~\ref{sec:grouping}): correlation grouping succeeds exactly when the
interacting features are mutually correlated and fails otherwise, and univariate
mutual information fails on jointly-informative features.
\end{itemize}

All code, synthetic generators, and scripts to reproduce every number in this
paper are released with the project.

\section{Related work}

\paragraph{Additive and glass-box models.} GAMs \citep{hastie1990gam} model the
log-odds as a sum of univariate shape functions. GA2M \citep{lou2013ga2m} adds a
controlled number of pairwise terms, and the Explainable Boosting Machine
\citep{caruana2015intelligible,nori2019interpretml} is its widely-used
gradient-boosted realization and our primary interpretable baseline. Neural
additive models \citep{agarwal2021nam} replace the shape functions with neural
networks but retain additivity in main (and optionally pairwise) effects. All of
these are, by design, limited to low-order interactions. PEAM departs by being
additive across \emph{groups} of features while placing an unconstrained learner
\emph{within} each group, trading per-feature shape plots for group-level reason
codes in exchange for arbitrary within-group interaction order.

\paragraph{Mixtures of experts.} PEAM's experts echo the mixture-of-experts
lineage \citep{jacobs1991moe}, but classical MoE uses a learned \emph{gate} to
select experts per input, which obscures attribution. PEAM instead assigns
experts to fixed feature groups and combines them additively, so that the
contribution of each expert is an exact, globally-consistent reason code. Recent
work on interpretable MoE \citep{ismail2023interpretable_moe} shares the goal of
transparent expert combination; our combiner is a transparent linear model in
log-odds space.

\paragraph{Log-odds attributions and faithfulness.} Working in log-odds (the
natural additive scale for a logistic model) connects PEAM to recent arguments
for additive evidence in interpretable modelling \citep{dervovic2022logodds}.
Unlike SHAP \citep{lundberg2017shap} and LIME \citep{ribeiro2016lime}, whose
attributions approximate a fixed black box and can disagree with it, PEAM's
reason codes are \emph{definitional}: they are the terms of the model's own
decision function and sum to it exactly.

\section{The Pattern-Expert Additive Model}
\label{sec:method}

\paragraph{Setup.} Let $x \in \RR^d$ with binary label $y \in \{0,1\}$ and base
rate $\pi = \Pr(y{=}1)$. Fix a partition of the feature indices into $K$ groups
$G_1,\dots,G_K$ (disjoint, covering all features). Write $x_{G_k}$ for the
sub-vector of features in group $k$, and let $\sigm$ be the logistic function and
$\logit = \sigm^{-1}$.

\paragraph{Experts.} For each group we fit a calibrated expert
$e_k : \RR^{|G_k|} \to (0,1)$ estimating $\Pr(y{=}1 \mid x_{G_k})$. In this paper
each expert is a histogram gradient-boosted tree trained on the training split
and calibrated by Platt scaling \citep{platt1999calibration} on a held-out
validation split, so that it reports a probability without leakage from the
calibration data into the base fit. Any model class can serve as an expert; the
only requirement is a calibrated probability.

\paragraph{Evidence and combiner.} Define the group evidence
\begin{equation}
z_k(x) \;=\; \logit e_k(x_{G_k}) \;-\; \logit \pi,
\end{equation}
i.e. how much group $k$ moves the log-odds away from the prior. PEAM combines
evidence linearly,
\begin{equation}
\logit\, p(x) \;=\; b_0 + \sum_{k=1}^K \beta_k\, z_k(x),
\label{eq:peam}
\end{equation}
with $(\beta_k, b_0)$ fit by logistic regression on the validation split
(the \emph{logistic} combiner). A parameter-free \emph{naive} variant sets
$\beta_k{=}1$ and $b_0 = \logit \pi$, recovering a naive-Bayes-style addition of
evidence.

\paragraph{Reason codes.} The per-group reason code is
\begin{equation}
r_k(x) \;=\; \beta_k\, z_k(x), \qquad \logit\, p(x) = b_0 + \sum_{k=1}^K r_k(x).
\end{equation}
This decomposition is \emph{exact by construction}: the reason codes are the
summands of the model's own decision function, not an approximation of it. A
positive $r_k$ is signed evidence from group $k$ for the positive class; the
global importance of a group is $\EE\,|r_k|$. Groups can be switched off
(drop $r_k$) for deployment in settings that lack a feature group, giving a
form of modular portability.

\begin{algorithm}[t]
\caption{PEAM training (logistic combiner)}
\label{alg:peam}
\begin{algorithmic}[1]
\Require train $(X^{tr},y^{tr})$, validation $(X^{val},y^{val})$, groups $\{G_k\}$
\State $\pi \gets \operatorname{mean}(y^{tr})$
\For{$k = 1,\dots,K$}
  \State fit base learner on $(X^{tr}_{G_k}, y^{tr})$
  \State calibrate it on $(X^{val}_{G_k}, y^{val})$ to obtain expert $e_k$
\EndFor
\State form evidence matrix $Z^{val}$ with $Z^{val}_{ik} = \logit e_k(x^{val}_{i,G_k}) - \logit\pi$
\State fit logistic regression of $y^{val}$ on $Z^{val}$ to obtain $(\beta, b_0)$
\State \Return $\{e_k\}, \beta, b_0, \pi$
\end{algorithmic}
\end{algorithm}

\paragraph{Properties.} (i) \emph{Exact additivity} of reason codes, as above.
(ii) \emph{Arbitrary within-group interaction order}: each $e_k$ may model any
interaction among the features of its group. (iii) \emph{Across-group additivity}:
Eq.~\eqref{eq:peam} assumes no interaction between groups --- this is PEAM's
inductive bias and, as we show, both its strength and its failure mode.
(iv) \emph{Calibration}: per-expert calibration plus a logistic combiner yields
well-calibrated probabilities (Section~\ref{sec:interp}).

\section{When does PEAM win?}
\label{sec:when}

Suppose the true conditional log-odds decomposes as
\begin{equation}
\logit \Pr(y{=}1\mid x) \;=\; c_0 + \sum_{k=1}^K f_k(x_{G_k}),
\label{eq:truth}
\end{equation}
i.e. it is additive \emph{across} the groups $\{G_k\}$ but each $f_k$ is an
arbitrary function of its group. Then:
\begin{itemize}
\item \textbf{PEAM is well-specified.} With a sufficiently flexible expert per
group, $e_k$ can match $f_k$ up to calibration, and the additive combiner in
Eq.~\eqref{eq:peam} matches the across-group structure. PEAM can represent
Eq.~\eqref{eq:truth} regardless of the \emph{order} of each $f_k$.
\item \textbf{GA2M / EBM is misspecified} whenever some $f_k$ contains a genuine
interaction of order $\geq 3$: its hypothesis class (main effects plus pairwise
terms) cannot express a third-order term, so it incurs irreducible approximation
error no matter how much data it sees.
\item \textbf{Linear models are misspecified} whenever any $f_k$ is nonlinear.
\end{itemize}
Conversely, if the truth contains interactions that \emph{cross} groups --- a term
depending jointly on features from $G_j$ and $G_k$ --- then PEAM's across-group
additivity (iii) is violated and PEAM is itself misspecified, while a black box
or even GA2M (through a cross-group pairwise term) may do better. Hence the
operative condition is not high interaction order alone but
\emph{high order \textbf{and} locality}: the interactions must live inside the
groups. We test both halves of this prediction below.

A second mechanism matters in practice. When many features are pure noise and
the informative features fall in one group, the logistic combiner drives the
noise groups' weights $\beta_k \to 0$, so PEAM performs \emph{implicit structured
feature selection}. This lets PEAM beat a monolithic black box that must find the
informative subset among all features --- a point borne out on MADELON.

\section{Controlled study: sweeping interaction order}
\label{sec:synthetic}

\paragraph{Generator.} We generate data in which the per-group signal is a
\emph{pure} order-$r$ interaction. For $K$ groups of $m$ features each, draw all
features i.i.d.\ $\mathcal N(0,1)$ (independent within and across groups) and set
\begin{equation}
z_k = \operatorname{standardize}\Big(\textstyle\prod_{j=1}^{r} x_{k,j}\Big), \qquad
\logit p = b + \sum_{k=1}^K w\, z_k, \qquad y \sim \text{Bernoulli}(\sigm(\logit p)),
\end{equation}
with $b$ solved to fix the positive rate (we use $10\%$). A product of $r$
independent mean-zero variables has zero correlation with any single feature and
with any lower-order product of its own factors, so the signal requires modelling
the \emph{full} order-$r$ interaction; the log-odds is additive across groups by
construction, so Eq.~\eqref{eq:truth} holds exactly. We add pure-noise features
to make grouping non-trivial.

\paragraph{Setup.} $K{=}4$ groups of $4$ features, $6$ noise features,
$10\%$ positives, stratified $60/20/20$ train/validation/test, three seeds. We
compare logistic regression, EBM (the interpretable baseline), PEAM with the true
grouping, and a black-box histogram GBM. Experts and the GBM use comparable
capacity.

\begin{table}[t]
\centering
\caption{Synthetic interaction-order sweep ($K{=}4$ groups, $10\%$ positives,
mean over 3 seeds). PEAM uses the true grouping. \texttt{gbm} is the
non-interpretable black-box reference; all others are interpretable.}
\label{tab:order}
\begin{tabular}{lrrrr}
\toprule
 & \multicolumn{4}{c}{ROC-AUC} \\
\cmidrule(lr){2-5}
interaction order & logreg & EBM & \textbf{PEAM} & gbm (black box) \\
\midrule
1 (linear)                 & 0.936 & 0.930 & 0.868 & 0.918 \\
2 (pairwise)               & 0.515 & 0.908 & 0.856 & 0.890 \\
3 (localized high-order)   & 0.485 & 0.668 & \textbf{0.795} & 0.731 \\
4 (localized high-order)   & 0.522 & 0.679 & \textbf{0.685} & 0.637 \\
\midrule
 & \multicolumn{4}{c}{PR-AUC} \\
\cmidrule(lr){2-5}
1 & 0.686 & 0.667 & 0.497 & 0.619 \\
2 & 0.121 & 0.618 & 0.472 & 0.555 \\
3 & 0.104 & 0.183 & \textbf{0.418} & 0.312 \\
4 & 0.128 & 0.179 & \textbf{0.254} & 0.166 \\
\bottomrule
\end{tabular}
\end{table}

\paragraph{Result (Table~\ref{tab:order}).} There is a clean crossover. At order
$1$ the truth is linear and logistic regression and EBM win; PEAM's flexible
experts are unnecessary overhead. At order $2$ the truth is exactly GA2M's design
point and EBM wins. At order $\geq 3$ EBM saturates --- it cannot represent the
interaction --- while PEAM with the true grouping becomes the \emph{single best
model, above the black box}, on both ROC-AUC and PR-AUC, because its inductive
bias (across-group additivity, free within-group order) matches the data exactly
whereas the black box must discover the group structure from limited data.

\paragraph{Robustness.} Is PEAM's edge over the black box a tuning or
small-sample artifact? Table~\ref{tab:nsweep} fixes order $3$ and gives the black
boxes every advantage: a strong early-stopped histogram GBM and XGBoost, across
sample sizes. PEAM leads at every $n$. The black box closes the gap as $n$ grows
(it eventually learns the block structure) but never overtakes the model that is
told it; EBM never catches up because order-3 terms lie outside its hypothesis
class.

\begin{table}[t]
\centering
\caption{Order-3 robustness (ROC-AUC, mean over 3 seeds). PEAM uses the true
grouping; the black boxes are strong and tuned (HistGBM: 2000 trees, early
stopping; XGBoost: 1500 trees, depth 6).}
\label{tab:nsweep}
\begin{tabular}{lrrrr}
\toprule
$n$ & \textbf{PEAM} & EBM & HistGBM (strong) & XGBoost \\
\midrule
2{,}000  & \textbf{0.651} & 0.636 & 0.603 & 0.622 \\
8{,}000  & \textbf{0.795} & 0.668 & 0.733 & 0.717 \\
30{,}000 & \textbf{0.884} & 0.725 & 0.856 & 0.845 \\
80{,}000 & \textbf{0.905} & 0.746 & 0.889 & 0.886 \\
\bottomrule
\end{tabular}
\end{table}

\section{Real-data study: where the regime occurs}
\label{sec:real}

We now ask whether the regime arises in public data. Table~\ref{tab:real}
summarizes four datasets chosen to span the possibilities; we describe the
grouping used for each.

\begin{table}[t]
\centering
\caption{Real-data hunt (ROC-AUC, mean over 3 seeds). \texttt{peam\_*} are
interpretable; \texttt{xgb\_*} is the black-box reference. Groupings are
domain-defined except where noted.}
\label{tab:real}
\begin{tabular}{llr}
\toprule
dataset & model & ROC-AUC \\
\midrule
\multirow{3}{*}{satimage (9 px $\times$ 4 bands)}
 & EBM & \textbf{0.961} \\
 & PEAM (semantic groups) & 0.948 \\
 & PEAM (correlation groups) & 0.941 \\
\midrule
\multirow{3}{*}{gas-drift (16 sensors $\times$ 8)}
 & logistic regression & 0.999 \\
 & EBM & 1.000 \\
 & PEAM (semantic groups) & 0.999 \\
\midrule
\multirow{4}{*}{HIGGS (per-object kinematics)}
 & EBM (low-level only) & \textbf{0.674} \\
 & PEAM (per-object groups) & 0.608 \\
 & XGBoost (low-level only) & 0.680 \\
 & EBM (with cross-object masses) & 0.796 \\
\midrule
\multirow{6}{*}{MADELON (5 informative + noise)}
 & EBM (full) & 0.800 \\
 & XGBoost (full) & 0.857 \\
 & \textbf{PEAM (correlation groups, full)} & \textbf{0.903} \\
 & EBM (shared feature selection) & 0.885 \\
 & XGBoost (shared feature selection) & 0.927 \\
 & \textbf{PEAM (shared feature selection)} & \textbf{0.912} \\
\bottomrule
\end{tabular}
\end{table}

\paragraph{MADELON: a localized high-order win.} MADELON
\citep{guyon2003madelon} has $500$ features, of which only $5$ are informative ---
placed on the vertices of a five-dimensional hypercube, so the class is defined
by a five-way interaction --- plus $15$ redundant features (random linear
combinations of the $5$) and $480$ noise probes. This is precisely a localized
high-order interaction embedded in noise. PEAM with correlation grouping
(Section~\ref{sec:grouping}) attains ROC-AUC $0.903$, beating EBM ($0.800$) and
full-feature XGBoost ($0.857$); it is the best interpretable model by a wide
margin and the best model overall \emph{except} for a black box equipped with a
separate feature-selection stage. To test fairness we gave every model the same
ExtraTrees-importance feature selection ($\approx 14$ features): even then PEAM
($0.912$) beats EBM ($0.885$), because selection removes noise but does not grant
EBM the fifth-order term it structurally lacks. Only the select-then-boost
black-box pipeline ($0.927$) edges PEAM --- at the cost of interpretability and an
extra stage. The win is stable: over five seeds PEAM averages $0.895$
(range $[0.863, 0.936]$) versus EBM $0.795$, Random Forest $0.768$, and
ExtraTrees $0.747$.

\paragraph{HIGGS: high-order but delocalized --- the contrast.} HIGGS
\citep{baldi2014higgs} is the canonical example of a dataset whose high-order
structure defeats additive models. Its $21$ low-level features are per-object
kinematics (lepton, missing energy, four jets); its $7$ high-level features are
\emph{invariant masses} that combine \emph{multiple} objects. Grouping by object
and applying PEAM to the low-level features gives ROC-AUC $0.608$, \emph{below}
EBM's $0.674$: the discriminative interactions cross object boundaries, violating
PEAM's across-group additivity. The point is sharpened by the last row: adding
the seven cross-object mass features lifts EBM from $0.674$ to $0.796$. HIGGS has
high-order signal, but it is \emph{delocalized}, so PEAM fails exactly as the
analysis in Section~\ref{sec:when} predicts. Locality, not interaction order
alone, is what PEAM requires.

\paragraph{satimage and gas-drift: not in the regime.} On satimage (remote
sensing; a $3{\times}3$ pixel neighbourhood over $4$ spectral bands) EBM already
reaches $0.961$ and PEAM with the genuine pixel/band grouping is slightly behind:
the signal is well captured by additive and pairwise terms. On gas-drift
($16$ chemical sensors $\times 8$ descriptors) every model, including logistic
regression, exceeds $0.99$ --- the task is near-separable, leaving no interaction
headroom. These are the common case for engineered tabular features, and they are
\emph{not} where PEAM helps.

\section{Which grouping works?}
\label{sec:grouping}

PEAM requires a grouping that places the interacting features together. We find
that the adequacy of a grouping is \emph{structure-dependent}, and that no single
automatic method is universal.
\begin{itemize}
\item \textbf{Correlation grouping} (hierarchical clustering on
$|\text{Spearman}|$) recovers the interacting set exactly when those features are
mutually correlated. On MADELON the $15$ redundant features are linear
combinations of the $5$ informative ones, so the useful set forms a correlated
cluster and correlation grouping isolates it from the $480$ noise probes --- which
is why PEAM wins there. In the synthetic generator, by contrast, the interacting
features are independent products with \emph{zero} pairwise correlation, so
correlation grouping cannot find them and PEAM with correlation (or random)
grouping collapses to near the additive baseline.
\item \textbf{Univariate mutual information} fails on jointly-informative
features: a pure high-order interaction has (near) zero marginal mutual
information with the label, so ranking features by univariate MI --- the classic
``MADELON trap'' --- selects the wrong features and PEAM gains nothing.
\end{itemize}
The practical consequence is that PEAM is most readily deployed when a grouping
is given by domain knowledge, and that \emph{interaction-aware} group discovery
(rather than correlation or marginal association) is the right open problem for
making PEAM automatic.

\section{Interpretability and calibration}
\label{sec:interp}

PEAM's reason codes are exact (Section~\ref{sec:method}); we also check that they
are \emph{faithful} and that the model is \emph{calibrated}. On a suite of
imbalanced public datasets, zeroing the evidence of the single most important
group at test time drops ROC-AUC by $\approx 0.14$ on average, confirming that
the reason codes identify genuine drivers rather than redundant signal. PEAM's
expected calibration error ($\approx 0.014$) is essentially tied with EBM's, the
best among the models we test, as a consequence of per-expert Platt calibration
followed by the logistic combiner \citep{platt1999calibration,guo2017calibration}.
Interpretability is at the \emph{group} level: PEAM reports which group drove a
decision and by how much, exactly, but --- unlike EBM's per-feature shape
functions --- it does not resolve attribution \emph{within} a group, whose expert
is a local black box. This is the representational price of capturing
within-group high-order interactions.

\section{Limitations}
\label{sec:limits}

We state the boundaries of the contribution plainly. (i) PEAM requires a grouping
that captures the interacting features; obtaining one automatically is unsolved
(Section~\ref{sec:grouping}). (ii) PEAM is dominated on linear, pairwise,
saturated, or cross-group-interaction data --- a large fraction of real tabular
problems. (iii) A black box combined with a separate feature-selection stage can
exceed PEAM even in its favourable regime (MADELON: $0.927$ vs $0.912$), trading
interpretability for a small accuracy gain. (iv) Interpretability is group-level,
not per-feature. PEAM is therefore best understood as a targeted instrument for a
specific, identifiable regime, not a replacement for GA2M or gradient boosting in
general.

\section{Conclusion}

We introduced PEAM, a group-additive classifier with flexible, calibrated
per-group experts and exact additive log-odds reason codes, and we characterized
--- honestly, including where it loses --- a previously uncharted point on the
accuracy--interpretability frontier. When the predictive signal consists of
high-order interactions localized within a known grouping, PEAM is the best
interpretable model and can beat strong black boxes; we demonstrated this in a
controlled order-sweep and on the MADELON benchmark, and we showed via HIGGS that
locality, not interaction order alone, is the operative condition. The natural
next step is interaction-aware group discovery, which would turn PEAM from a
model that \emph{exploits} a known grouping into one that \emph{finds} the regime
on its own.

\paragraph{Reproducibility.} The PEAM implementation, the synthetic generator,
and scripts that regenerate every table (\texttt{run\_interaction\_locality.py},
\texttt{run\_real\_data\_hunt.py}) are released with the project. All real
datasets are public and fetched without credentials. Experiments run on CPU.

\small
\bibliographystyle{plainnat}
\bibliography{refs}

\end{document}
